% Part: incompleteness
% Chapter: incompleteness-provability
% Section: tarski-thm

\documentclass[../../../include/open-logic-section]{subfiles}

\begin{document}

\olfileid[fr]{inc}{inp}{tar}

\olsection{L'indéfinissabilité de la vérité}

La notion de \emph{définissabilité} suppose que
l'on ait donné une sémantique formelle du langage
de l'arithmétique. Nous avons décrit un ensemble de
!!{formula}s et d'!!{sentence}s de ce langage.
Selon l'« interprétation voulue », ces
!!{sentence}s expriment des assertions sur les
entiers naturels; une telle assertion peut être
vraie ou fausse. Soit $\Struct{N}$ la
!!{structure} de domaine~$\Nat$ qui donne aux
symboles du langage de l'arithmétique leur
interprétation usuelle. Alors $\Sat{N}{!A}$
signifie que « $!A$ est vraie dans l'interprétation
usuelle ».

\begin{defn}
Une relation $R(x_1,\dots,x_k)$ sur les entiers
naturels est \emph{définissable} dans~$\Struct{N}$
si et seulement s'il existe une !!{formula}
$!A(x_1,\dots,x_k)$ du langage de l'arithmétique
telle que, pour tous $n_1,\dots,n_k$,
$R(n_1,\dots,n_k)$ si et seulement si
$\Sat{N}{!A(\num n_1,\dots,\num n_k)}$.
\end{defn}

Autrement dit, une relation est définissable dans
$\Struct{N}$ si et seulement si elle est
représentable dans la théorie $\Th{TA}$, où
$\Th{TA}=\Setabs{!A}{\Sat{N}{!A}}$ est l'ensemble
des !!{sentence}s vrais de l'arithmétique.
Si cette équivalence ne vous paraît pas immédiate,
reprenez les définitions pour vous en convaincre.

\begin{lem}
Toute relation calculable est définissable
dans~$\Struct{N}$.
\end{lem}

\begin{proof}
On vérifie facilement que la !!{formula}
représentant une relation dans~$\Th{Q}$ définit
la même relation dans~$\Struct{N}$.
\end{proof}

La réciproque est-elle vraie? Toute relation
définissable dans~$\Struct{N}$ est-elle calculable?
Non. En voici un exemple.

\begin{lem}
La relation d'arrêt est définissable
dans~$\Struct{N}$.
\end{lem}

\begin{proof}
Rappelons que le théorème de la forme normale de
Kleene affirme que toute fonction partiellement
calculable~$f$ possède un indice~$e$ tel que,
pour tout $x\in\Nat$,
$f(x)=U(\umin{s}{T(e,x,s)})$, où $U$ et $T$
sont récursives primitives : leurs évaluations
terminent donc toujours.
Ainsi, $f(x)$ est définie, c'est-à-dire que le
calcul s'arrête, si et seulement s'il existe
un~$s$ tel que $T(e,x,s)$.

Soit maintenant $H$ la relation d'arrêt :
\[
H=\Setabs{\tuple{e,x}}{\lexists[s][T(e,x,s)]}.
\]
Soit $!D_T$ une !!{formula} qui définit~$T$
dans~$\Struct{N}$. Alors
\[
H=\Setabs{\tuple{e,x}}{
\Sat{N}{\lexists[s][!D_T(\num e,\num x,s)]}},
\]
si bien que
$\lexists[s][!D_T(z,x,s)]$ définit~$H$
dans~$\Struct{N}$.
\end{proof}

\begin{prob}
Montrez que la relation
$Q(n)\defiff n\in
\Setabs{\Gn{!A}}{\Th{Q}\Proves !A}$
est définissable en arithmétique.
\end{prob}

Qu'en est-il de $\Th{TA}$ elle-même?
Est-elle définissable en arithmétique?
Autrement dit, l'ensemble
$\Setabs{\Gn{!A}}{\Sat{N}{!A}}$
est-il définissable en arithmétique?
Le théorème de Tarski répond par la négative.

\begin{thm}
\ollabel{thm:tarski}
L'ensemble des !!{sentence}s vrais de
l'arithmétique n'est pas définissable en
arithmétique.
\end{thm}

\begin{proof}
Supposons que $!D(x)$ définisse cet ensemble :
$\Sat{N}{!A}$ si et seulement si
$\Sat{N}{!D(\gn{!A})}$. Par le lemme du point
fixe, il existe un !!{sentence}~$!A$ tel que
$\Th{Q}\Proves !A\liff\lnot !D(\gn{!A})$,
et donc
$\Sat{N}{!A\liff\lnot !D(\gn{!A})}$.
Ainsi, $\Sat{N}{!A}$ si et seulement si
$\Sat{N}{\lnot !D(\gn{!A})}$, ce qui contredit
le fait que $!D(y)$ définisse l'ensemble des
!!{sentence}s vrais de l'arithmétique.
\end{proof}

Tarski a appliqué cette analyse à une notion
philosophique plus générale de vérité. Pour tout
langage~$L$, une notion de vérité adéquate devrait,
selon lui, satisfaire, pour chaque énoncé~$X$, le
schéma suivant :
\begin{quote}
« $X$ » est vrai si et seulement si $X$.
\end{quote}
Son exemple souvent cité, formulé en français, est :
\begin{quote}
L'énoncé « La neige est blanche » est vrai si et
seulement si la neige est blanche.
\end{quote}
Cependant, dans tout langage assez puissant pour
représenter la fonction diagonale, et pour tout
prédicat linguistique~$T(x)$, on peut construire
un !!{sentence}~$X$ qui satisfait :
« $X$ si et seulement si non
$T(\text{« $X$ »})$ ».
Puisque nous ne voulons pas qu'un prédicat de
vérité déclare certains !!{sentence}s à la fois
vrais et faux, Tarski en conclut qu'on ne peut
définir un prédicat de vérité pour tous les
!!{sentence}s d'un tel langage sans sortir, d'une
façon ou d'une autre, du langage lui-même.
Autrement dit, le prédicat de vérité d'un tel langage
ne peut pas être défini dans ce langage.

\end{document}
